% Seccion 7 de la rubrica: discusion y conclusiones (1 a 2 paginas), con los
% resultados mas relevantes, recomendaciones y las preguntas nuevas que deja
% el analisis.

\section{Discusión y conclusiones}
\label{sec:conclusiones}

\subsection{Discusión}

El objetivo era clasificar cada día-estación como excedencia o no del MDA8
a 51 ppb a partir de las condiciones del mismo día, con AUC de al menos
0.75, e identificar las condiciones latentes que explican esa
clasificación. Se cumple: los tres clasificadores superan el criterio, y las
técnicas de interdependencia producen tres ejes con lectura física estable
fuera del periodo de ajuste.

La excedencia de ozono en el AMM es, ante todo, un fenómeno estacional y
regional. La temporada es el predictor de mayor efecto en los modelos
día-estación ---la primavera triplica los momios del invierno--- y el
forzamiento fotoquímico, el eje de temperatura, radiación y sequedad, es el
factor continuo dominante. Que el forzamiento pese más que la carga de
precursores es coherente con la química del ozono como contaminante
secundario: los ingredientes son necesarios, pero es la energía disponible
la que fija la velocidad de las reacciones. La ventilación es el único
mecanismo protector medido, y el AF muestra que es el viento, no la presión,
lo que define ese eje.

Sobre ese fondo estacional hay una señal espacial que la meteorología medida
no explica. A igual temporada y condiciones, el Noroeste y el Suroeste casi
cuadruplican los momios del Noreste, y la tasa de excedencia de una estación
correlaciona 0.83 con su elevación. El viento predominante del oriente
arrastra los precursores hacia el pie de la sierra, donde la orografía frena
su salida; el ordenamiento poniente--oriente es estable en las cuatro
temporadas. Los contrastes de zona descansan, sin embargo, en 15 unidades
espaciales, y su precisión es la que 15 unidades pueden dar.

La dependencia temporal es el tercer hallazgo. La autocorrelación de un día
para otro sobrevive a la descomposición estacional en las cuatro series de
área metropolitana, y un solo rezago de la respuesta captura casi toda esa
dependencia: haber excedido ayer casi cuadruplica los momios de exceder hoy,
controlando por factores y calendario. El modelo autorregresivo distingue
días dentro de una temporada, lo que la climatología no puede hacer.

Dos limitaciones acotan todo lo anterior. La primera es el efecto de año:
2024 y 2025 tienen un nivel de excedencia que los factores no explican, y
ningún modelo generaliza igual de bien a 2021 que a 2024. Con estos datos
no puede decidirse si es un cambio de emisiones, de instrumentación o dos
años meteorológicamente atípicos. La segunda es la ausencia de compuestos
orgánicos volátiles: sin ellos, los ejes se describen como consistentes con
combustión, ventilación o fotoquímica, y ninguna afirmación sobre el régimen
que limita la formación de ozono es contestable.

\subsection{Conclusiones}

\begin{enumerate}[label=\roman*.]
  \item Los tres clasificadores alcanzan AUC de 0.846 a 0.898 en 2025 y
    superan el criterio de logro. La logística y el discriminante son
    equivalentes en la práctica, de modo que la interpretación de las razones
    de momios no depende del supuesto de normalidad.
  \item La excedencia es principalmente estacional: la temporada pesa más
    que cualquier factor continuo, y entre estos domina el forzamiento
    fotoquímico (2.78 por desviación estándar) sobre la combustión primaria
    (1.92); la ventilación protege (0.90).
  \item Existe una señal espacial no explicada por la meteorología medida:
    el poniente del AMM casi cuadruplica los momios del Noreste, en
    correspondencia con la elevación y el viento predominante.
  \item El calendario laboral apenas influye; solo el domingo tiene un
    efecto claro y pequeño, y el del sábado desaparece al corregir los
    errores estándar por dependencia.
  \item La dependencia día a día es genuina y un modelo AR(1) la aprovecha:
    supera a la climatología y a la persistencia con generalización esperada
    de 0.878.
\end{enumerate}

\subsection{Recomendaciones}

Para la gestión de la calidad del aire, el modelo autorregresivo sobre la
serie de área metropolitana puede operar como sistema de \emph{nowcasting}
de excedencia con las mediciones del propio día, con un umbral de decisión
que privilegie la sensibilidad, porque el costo de no avisar de un día malo
supera al de una alerta de más. Las zonas Noroeste y Suroeste merecen
vigilancia y medidas de control específicas, ya que su riesgo excede lo que
explican sus condiciones meteorológicas. Y la incorporación de mediciones de
compuestos orgánicos volátiles a la red del SIMA es la única vía para pasar
de asociaciones a mecanismo.

Para el análisis, tres mejoras metodológicas quedan identificadas: estimar
el parámetro de Box-Cox solo sobre el ajuste, para cerrar la única fuga que
persiste; incluir un término de media móvil sobre los residuales, que el
residuo de segundo orden del modelo autorregresivo justifica; y extender el
modelo a pronóstico rezagando los factores, con validación propia.

\subsection{Preguntas nuevas}

\begin{itemize}
  \item ¿Qué produjo el escalón de 2024? Distinguir entre cambio de
    emisiones, de instrumentación y años atípicos exige inventarios de
    emisiones y bitácoras de calibración que este trabajo no tuvo.
  \item ¿Qué característica del poniente explica el exceso de riesgo que
    la zona conserva tras controlar por meteorología? Uso de suelo,
    distancia a fuentes industriales y recirculación orográfica son las
    candidatas.
  \item ¿Cuánto del efecto de la temporada es fotoquímica y cuánto es
    régimen de vientos? Los factores lo capturan en parte, pero la temporada
    conserva un efecto propio que sugiere una variable estacional no medida.
  \item ¿Hasta qué horizonte es útil el pronóstico? El tiempo de mezcla de
    1.3 días de la cadena binaria sugiere que más allá de dos o tres días la
    predicción converge a la climatología del mes.
\end{itemize}
